\section{Results}\label{sec:results}

\subsection{Scores and thresholds}
CD$_5$ is defined for 86,904 of the 86,950 scored papers; the other 46 (25 controls, two papers retracted inside their
window and 19 papers of the reference sample) have no paper in the window that cites them or one of their references,
so $n_i+n_j+n_k=0$. Thus, 52,122 of the full pool of 52,147 controls have a defined score; the 25 undefined scores
are excluded, not treated as zero.
The distribution is concentrated near zero: in the scored set the 5th, 50th and 95th
percentiles are $-0.0218$, $-0.0010$ and $0.0059$. The large $n_k$ term compresses the scores
(\ref{sec:appendix}). Papers without a citing paper in the window but with $n_k>0$ score exactly zero. CD$_5 = 1$ is
rare (0.04\% of defined scores) because papers without references are excluded. In the reference sample, the yearly
top-5\% threshold falls from 0.021 (2000) to between 0.005 and 0.006 (2014--2017), and the top-1\% threshold from 0.13
to between 0.04 and 0.06. The mean CD$_5$ per publication year stays between $-0.0049$ and $0.0009$.

\subsection{RQ1: retracted papers versus matched controls}
The comparison population is shown in Fig.~\ref{fig:population}. Control percentages below use its defined-score
comparison group, with each stratum weighted by its number of scored primary-set papers, as for the mean difference.
The stratum-weighted mean CD$_5$ is $-0.002165$ for retracted papers and $-0.002363$ for their controls, a difference of
$0.00020$ (95\% CI $-0.00061$ to $0.00109$).
The post hoc 90\% interval, $-0.00051$ to $0.00093$, rules out absolute
mean differences larger than 0.00093, about 28\% of the difference between reviews and other papers (below).
Retracted papers have a positive score somewhat
more often within the comparison strata (923 of 4,586, 20.1\%, against 17.3\% of their controls, 95\% bootstrap CI
16.0--18.6; stratified odds ratio 1.23, 1.13--1.34).

\subsection{RQ2: reason classes}
The mean differences for misconduct (2,742 compared papers) and paper-mill production (276) are small
and their 95\% intervals include zero (Table~\ref{tab:main}). The post hoc 90\% intervals are $-0.00010$ to $0.00074$
and $0.00001$ to $0.00111$, respectively.
Their odds of reaching the top 5\% are close to one but imprecise (Table~\ref{tab:main},
Fig.~\ref{fig:forest}). None of the 276 paper-mill papers reaches the top 1\%. Papers retracted for honest error (294)
score lower than their controls (difference $-0.0041$, $-0.0077$ to $-0.0007$), that is, they are slightly more
consolidating.

\subsection{RQ3: the most disruptive papers}
In the pre-specified five-year analysis, retracted papers are not over-represented among the most disruptive papers of their year
(Table~\ref{tab:main}). 3.1\% of them (2.6--3.6) reach the top-5\% threshold, against 3.1\% (2.8--3.5) of their
controls, weighted in the same way; within strata the odds ratio is 0.99 (0.83--1.20). For the top 1\%, 17
retracted papers (0.37\%, 0.23--0.59) reach the threshold, against 0.54\% (0.43--0.66) of their controls; the odds
ratio of 0.67 (0.40--1.12) is imprecise. Both groups reach the yearly thresholds less often than
the nominal 5\% and 1\%: they come from the journals in which retractions occur, whereas the thresholds come from the
reference sample of all eligible papers.

\begin{table}[t]
\centering
\caption{CD$_5$ of retracted papers (primary set) versus controls from the same journal and year, by reason class.
Diff.: stratified mean difference ($\times 10^{-3}$, 95\% bootstrap CI). OR: Mantel--Haenszel odds ratio (95\% CI).
$n$: retracted papers in strata with a defined control score, not all scored retracted papers.}
\label{tab:main}
\resizebox{\linewidth}{!}{%
\begin{tabular}{lrrrrr}
\toprule
 & $n$ & Diff. ($\times10^{-3}$) & OR, CD$_5>0$ & OR, top 5\% & Top 1\% ($n$) \\
\midrule
All retracted & 4,586 & 0.20 ($-0.61$, 1.09) & 1.23 (1.13, 1.34) & 0.99 (0.83, 1.20) & 17 \\
Misconduct & 2,742 & 0.32 ($-0.21$, 0.82) & 1.29 (1.15, 1.44) & 0.93 (0.71, 1.22) & 5 \\
Paper mill & 276 & 0.58 ($-0.14$, 1.23) & 1.20 (0.93, 1.57) & 1.10 (0.53, 2.29) & 0 \\
Honest error & 294 & $-4.08$ ($-7.68$, $-0.73$) & 1.15 (0.82, 1.60) & 1.02 (0.55, 1.88) & 2 \\
Other & 1,274 & 0.84 ($-1.72$, 3.97) & 1.18 (1.02, 1.36) & 1.06 (0.79, 1.41) & 10 \\
\bottomrule
\end{tabular}}
\end{table}

\subsection{Robustness}
Table~\ref{tab:rob} gives the pre-specified checks. Excluding CD$_5 = 1$ papers (a), reviews (b) or using the window
$t+1$ to $t+5$ (d) leaves the results unchanged. Three checks move them. Splitting strata by reference-count decile
(c) gives a positive mean difference and an odds ratio of 1.47 (1.15--1.88) for the top 5\%; this check keeps only the
3,664 retracted papers that have a control in the same decile. The reference-length analysis below describes an
imbalance in the direction of this shift. Restricting to papers with at least five citations in the window (f) gives a
small positive mean difference. Including the 5,863 papers retracted inside their window (e) raises the odds ratio for
a positive score from 1.23 to 1.50 (1.42--1.58), while the top-5\% odds ratio stays at 1.00.

\begin{table}[t]
\centering
\caption{Robustness checks for the primary set (CD$_5$) and the post hoc three-year window (CD$_3$).
Diff.\ in units of $10^{-3}$; all estimates have 95\% intervals in parentheses.}
\label{tab:rob}
\resizebox{\linewidth}{!}{%
\begin{tabular}{lrrr}
\toprule
\multicolumn{4}{l}{CD$_5$, 2000--2017} \\
Check & Diff. & OR, CD$>0$ & OR, top 5\% \\
\midrule
Main analysis & 0.20 ($-0.61$, 1.09) & 1.23 (1.13, 1.34) & 0.99 (0.83, 1.20) \\
(a) no CD $=1$ & 0.05 ($-0.50$, 0.59) & 1.23 (1.13, 1.33) & 0.99 (0.82, 1.19) \\
(b) no reviews & 0.14 ($-0.74$, 1.18) & 1.23 (1.12, 1.34) & 0.98 (0.81, 1.19) \\
(c) + reference deciles & 1.23 (0.39, 2.23) & 1.38 (1.24, 1.53) & 1.47 (1.15, 1.88) \\
(d) window from $t+1$ & 0.31 ($-0.63$, 1.35) & 1.24 (1.14, 1.35) & 1.01 (0.84, 1.21) \\
(e) all retracted & 1.44 (0.78, 2.21) & 1.50 (1.42, 1.58) & 1.00 (0.90, 1.12) \\
(f) $\ge 5$ citations & 0.62 (0.07, 1.15) & 1.28 (1.17, 1.41) & 1.09 (0.88, 1.35) \\
\midrule
\multicolumn{4}{l}{CD$_3$, 2000--2019 (post hoc)} \\
Check & Diff. & OR, CD$>0$ & OR, top 5\% \\
\midrule
Main analysis & 0.39 ($-0.27$, 1.14) & 1.30 (1.22, 1.39) & 1.12 (0.98, 1.28) \\
(a) no CD $=1$ & 0.32 ($-0.17$, 0.83) & 1.30 (1.22, 1.39) & 1.12 (0.98, 1.28) \\
(b) no reviews & 0.39 ($-0.32$, 1.17) & 1.31 (1.22, 1.40) & 1.13 (0.98, 1.30) \\
(c) + reference deciles & 1.26 (0.52, 2.02) & 1.33 (1.23, 1.44) & 1.46 (1.23, 1.73) \\
(d) window from $t+1$ & 0.54 ($-0.22$, 1.38) & 1.32 (1.24, 1.40) & 1.15 (1.01, 1.31) \\
(e) all retracted & 1.27 (0.67, 1.90) & 1.54 (1.47, 1.61) & 1.16 (1.06, 1.26) \\
(f) $\ge 5$ citations & 0.78 (0.23, 1.36) & 1.39 (1.29, 1.49) & 1.32 (1.12, 1.56) \\
\bottomrule
\end{tabular}}
\end{table}

\subsection{Exploratory: why the sign shift?}
\paragraph{A review positive control}
Reviews among the never-retracted controls score lower than the other controls of the same journal and year: the mean
difference in CD$_5$ is $-0.0033$ ($-0.0066$ to $-0.0011$), and reviews reach the top 5\% (odds ratio 0.58,
0.47--0.72) and the top 1\% (0.39, 0.21--0.73) much less often (6,032 reviews, 4,832 of them in the 2,083 strata that
also contain other controls; Fig.~\ref{fig:forest}). The design therefore detects a contrast expected from the
construction of the index, with sample sizes similar to those of the main analysis.
The odds ratio for a positive CD$_5$ is 0.95 (0.86--1.04), and the mean DI$_{nok}$ difference is
$-0.076$ ($-0.195$ to $0.033$; 4,736 compared reviews). Both intervals include the null: the positive control
detects the CD$_5$ mean and top-tier contrasts, but not a clear sign or DI$_{nok}$ mean contrast.
The bootstrap resamples weighted journal--year contrasts, whose values and weights differ between the review and
retraction comparisons; similar numbers of papers therefore do not imply similar interval widths.

\paragraph{Citation composition}
The variant without $n_k$ reveals a citation-composition difference that CD$_5$ compresses.
After excluding notice, Comment/Letter and retracted citing records, the mean DI$_{nok}$ gap is 0.082
(95\% journal-bootstrap CI 0.049--0.110; Table~\ref{tab:citers}).
The matched means are $-0.2995$ for retracted papers and $-0.3811$ for controls:
on average, 65.0\% of papers citing a retracted paper also cite a reference, against 69.1\% for controls.
The comparison retains 4,504 retracted papers and 30,389 controls in 2,190 strata and 972 journals
(Fig.~\ref{fig:population}).

Excluding these citing records increases rather than removes the gap. With the same journal-bootstrap design,
the unfiltered gap is 0.067 (0.034--0.095), becomes 0.069 after notice removal and 0.070 after also removing
comments and letters. The full filter makes 13 previously defined primary scores undefined; comparison strata
are re-evaluated, not filled with zeros.
Adjusting the DI$_{nok}$ mean for log reference and remaining citing-paper counts within strata gives a gap
of 0.097 (0.065--0.127). Excluding the ten journals with most original primary-set retractions reduces the filtered
gap to 0.053 (0.033--0.074); with both covariates it is 0.070 (0.049--0.091).
Class-specific and population diagnostics are in \ref{sec:appendix}.

CD and DI$_{nok}$ have the same sign. In the original counts, excluding uncited papers leaves a positive-score
odds ratio of 1.20 (1.11--1.31); holding reference and citation counts fixed gives 1.32 (1.21--1.43).
These sign models and the unfiltered intervals in Fig.~\ref{fig:forest} are retained separately in
Table~\ref{tab:posthoc}. DI$_{nok}$ cannot define a top tier: its yearly 95th and 99th percentiles equal 1.

\begin{table}[!htbp]
\centering
\caption{Exploratory citing-record sensitivity of the DI$_{nok}$ mean. All intervals are 95\% whole-journal
bootstrap intervals; original journal--year matching is retained. Filters remove citing records, not focal references.
Adjusted rows are weighted within-stratum linear models. $n$: compared retracted papers.}
\label{tab:citers}
\small
\begin{tabular}{@{}>{\raggedright\arraybackslash}p{0.48\linewidth}lr@{}}
\toprule
Analysis & Mean gap (95\% CI) & $n$ \\
\midrule
CD$_5$ window: unfiltered & 0.067 (0.034, 0.095) & 4,517 \\
Exclude notices & 0.069 (0.037, 0.096) & 4,517 \\
Also exclude Comment/Letter & 0.070 (0.037, 0.098) & 4,513 \\
Also exclude retracted citing papers & 0.082 (0.049, 0.110) & 4,504 \\
Full filter, adjust log references & 0.094 (0.060, 0.124) & 4,504 \\
Full filter, also adjust log citations & 0.097 (0.065, 0.127) & 4,504 \\
Full filter, exclude top ten journals & 0.053 (0.033, 0.074) & 3,013 \\
Same exclusion, both covariates & 0.070 (0.049, 0.091) & 3,013 \\
CD$_3$ window: full filter & 0.090 (0.064, 0.114) & 7,088 \\
CD$_3$: also exclude top ten journals & 0.082 (0.064, 0.099) & 5,057 \\
\bottomrule
\end{tabular}
\end{table}

\subsection{Exploratory: reference length and the top tier}
The top-tier comparison is sensitive to reference length. In the conditional logistic regression, the odds ratio for
reaching the top 5\% of CD$_5$ is 0.99 without covariates and 1.42 (1.16--1.75) with $\log(1+\text{references})$;
adding $\log(1+\text{citing papers in the window})$ gives 1.19 (0.96--1.47). Short reference lists favour high scores
(coefficient $-1.64$ per unit of log references). The matched diagnostic, using the same log transformation
as this model, gives means of $\log(1+\text{references})$ of 3.6418 for the 4,586 compared retracted papers and 3.5575
for their 31,662 controls, weighted within strata. The difference is $+0.084$ (95\% strata-bootstrap CI
0.058--0.109). This imbalance pushes retracted papers' unadjusted scores down.
The corresponding matched difference in mean $\log(1+\text{citing papers})$ is $+0.311$ (0.280--0.341).

Restricting both groups by PubMed publication types retains the reference-length imbalance and reference-adjusted
top-tier excess. The filter, full estimates and population accounting are in \ref{sec:appendix}.

\subsection{Exploratory: what happens around the notice?}
Removing editorial notice records changes the primary positive-score estimate little, but reduces it for papers
retracted inside the scoring window (\ref{sec:appendix}).
Beyond notices, the citing papers change around the retraction. For 5,100 retracted papers
from 2000--2019 with citing papers before and after the notice, DI$_{nok}$ rose by 0.174 from the three years before
the notice year to the three years after it, against 0.100 for their controls over the same calendar years: a
difference in differences of 0.074 (0.050--0.097). Their citations fell from 3.9 to 2.6 per year, while those of the
controls rose from 3.1 to 4.7. The change began before the notice, as the placebo check confirms
(\ref{sec:appendix}).
In a separate balanced population of 3,083 papers, Figure~\ref{fig:event} plots the mean paired annual gap
in DI$_{nok}$: 0.053 (95\% CI 0.020--0.083) three years before the notice, 0.092 (0.059--0.124) in the notice year,
and 0.201 (0.162--0.242) two years after it. These estimates use 2,220, 2,440 and 1,982 defined pairs,
respectively; uncited annual scores are undefined.

\subsection{Exploratory: a shorter scoring window}
The shorter window admits more recent papers, including more paper-mill papers (population counts in \ref{sec:appendix}).
The mean difference remains small,
$0.00039$ ($-0.00027$ to $0.00114$), and the odds ratio for a positive score is 1.30
(1.22--1.39); top 5\% 1.12 (0.98--1.28); top 1\%
0.89 (0.63--1.28). Unlike the pre-specified five-year comparison, these post hoc top-5\% estimates lean toward a
modest excess: the odds ratios are 1.15 (1.01--1.31) when the window starts at $t+1$, 1.16 (1.06--1.26) for all
retracted papers, and 1.32 (1.12--1.56) among papers with at least five window citations (Table~\ref{tab:rob}).
For paper-mill papers the mean difference is $0.00039$ ($-0.00009$ to $0.00082$), the top-5\% odds
ratio 1.23 (0.83--1.81), and none of the 708 reaches the top 1\% (Fig.~\ref{fig:forest}).
The article-type restriction strengthens the apparent three-year top-tier excess, whereas author-name exclusions
attenuate it. Full estimates are in \ref{sec:appendix}; the excess is not clear under every sensitivity restriction.
The fully filtered DI$_{nok}$ mean gap remains positive in this window, including after the top-ten-journal
exclusion (Table~\ref{tab:citers}).
